
Para la validación de este sistema, se ha hecho uso de un modelo matemático que simula el comportamiento de una neurona, este es el modelo de Hindmarsh-Rose, que es capaz de generar ráfagas de potenciales o \textit{spikes} tanto de duración regular como irregular.  El modelo consiste del siguiente conjunto de ecuaciones diferenciales no lineales acopladas, como se puede ver a continuación.

\begin{equation} {Sistema de ecuaciones diferenciales del modelo neuronal de Hindmarsh-Rose.}
\begin{aligned}
\frac{dx(t)}{dt} &= y(t) + 3x^2(t) - x^3(t) - z(t) + e \\
\frac{dy(t)}{dt} &= 1 - 5x^2(t) - y(t) \\
\frac{1}{\mu} \frac{dz(t)}{dt} &= - z(t) + S[x(t) + 1.6]
\end{aligned}
\end{equation}


La variable $x$ representa el potencial de membrana de la neurona, las variables $y$ y $z$ representan la acción combinada de canales iónicos, siendo $y$ de acción rápida y $z$ de acción lenta. En la primera ecuación $e$ es un estímulo externo a la neurona, $\mu$ es la constante de tiempo de la dinámica lenta y $S$ es un parámetro que controla la influencia de la variable rápida $x$ en la lenta $z$ \cite{Hindmarsh1984}. El valor de $e$, hace que el modelo siga un patrón regular (Figura \ref{FIG:HRREG}) o que entre en modo caótico (Figura \ref{FIG:HRCHAO})


\begin{figure}[Ejecución de Hindmarsh-Rose Regular]{FIG:HRREG}{Ejecución del modelo de Hindmarsh-Rose con una $e <= 3$ con el programa \textit{Pico Neuron} (Sección \ref{SEC:PICONEURON}).}
    \image{0.9\textwidth}{}{hindmarsh-rose}
\end{figure}


\begin{figure}[Ejecución de Hindmarsh-Rose Caótico]{FIG:HRCHAO}{Ejecución del modelo de Hindmarsh-Rose con una $e > 3$ con el programa \textit{Pico Neuron} (Sección \ref{SEC:PICONEURON}).}
    \image{0.9\textwidth}{}{hindmarsh-rose-chaotic}
\end{figure}